\chapter[RAPPORTS D'OBJETS ET D'ITEMS]{Rapports d'Objets et d'Items}\label{rapports-dobjets-et-ditems}

\newcommand{\chElevenHead}[1]{\shortstack{\sffamily\bfseries\itshape #1}}
\newcommand{\chElevenDataExtHead}{%
  \makebox[\linewidth][c]{%
  \begin{tabular}{@{}>{\centering\arraybackslash\sffamily\bfseries\itshape}p{\dimexpr\linewidth+2\tabcolsep\relax}@{}}
  Course/Speed \\
  \hline
  Power/Height/Gain/Dir \\
  \hline
  Radio Range \\
  \hline
  DF Signal Strength \\
  \hline
  Area Object
  \end{tabular}}%
}
\newcommand{\chElevenBytes}[1]{%
  \makebox[0pt][r]{Bytes:\hspace{1.2em}}%
  \makebox[\linewidth][c]{#1}%
}
\newcommand{\chElevenExampleText}[1]{{\rmfamily\small #1}}

\begin{aprssideblock}{Objets et Items}
Toute station APRS peut reporter manuellement la position d'une entité
APRS (par exemple une autre station ou un phénomène météorologique). Cela
est prévu pour les cas où l'entité n'est pas capable de reporter sa propre
position.

APRS fournit deux types de rapport pour cela :

\begin{itemize}
\tightlist
\item
  \textbf{Object Reports}
\item
  \textbf{Item Reports}
\end{itemize}

Les Object Reports spécifient la position d'un Object, peuvent avoir un
timestamp optionnel, et peuvent inclure des informations de cap/vitesse ou
d'autres données étendues. Les Object Reports sont prévus principalement
pour tracer la position d'objets mobiles (par exemple les engins spatiaux,
les orages, les coureurs de marathon sans tracker).

Les Item Reports spécifient la position d'un Item, mais ne peuvent pas
avoir de timestamp. Bien que les Item Reports puissent aussi inclure des
données de cap/vitesse ou d'autres données étendues, ils sont réellement
prévus pour des choses inanimées qui sont occasionnellement postées sur
une carte (par exemple les points de contrôle de marathon ou les postes de
secours). Sinon, ils sont traités de la même manière que les Object
Reports.

Les Objects se distinguent les uns des autres par des noms d'Object
différents. De même, les Items se distinguent les uns des autres par des
noms d'Item différents.

\textbf{Recommandation d'implémentation :} lorsqu'un Object/Item APRS est
affiché à l'écran, l'indicatif de la station qui envoie le rapport doit
être associé à l'Object/Item.
\end{aprssideblock}

\begin{aprssideblock}{Remplacer un\\Object / Item}
Un précepte fondamental d'APRS est que n'importe quelle station peut
reprendre la responsabilité du report d'un Object ou d'un Item APRS, en
transmettant simplement un nouveau rapport portant le même nom
d'Object/Item.

Le rapport de remplacement peut spécifier l'emplacement existant ou un
nouvel emplacement. La station d'origine cessera de transmettre le rapport
d'Object/Item lorsqu'elle verra un rapport entrant portant le même nom
provenant d'une autre station.
\end{aprssideblock}

\begin{aprssideblock}{Tuer un\\Object / Item}
Pour tuer un Object/Item, une station transmet un nouveau rapport
d'Object/Item, avec un caractère de « kill » suivant le nom de
l'Object/Item.

\textbf{Recommandation d'implémentation :} lorsqu'un Object/Item est tué,
il doit être retiré de l'affichage à l'écran. Cependant, les données
associées à l'Object/Item doivent être conservées en interne au cas où
elles seraient nécessaires ultérieurement.
\end{aprssideblock}

\par\noindent{\sffamily\bfseries Additional Amplifying Comments about Objects}\par

\begin{aprssideblock}{Format Object\\Report}
Un Object Report possède un champ de nom d'Object fixe de 9 caractères,
qui peut être constitué de n'importe quels caractères ASCII imprimables, y
compris des espaces incorporés. Des espaces de fin sont utilisés pour
rendre le champ large de 9 caractères. Ceux-ci ne font pas partie du nom.
Les noms d'Object sont sensibles à la casse.

Le \lit{;} est l'APRS Data Type Identifier d'un Object Report, et un
\lit{*} ou un \lit{\_} sépare le nom d'Object du reste du rapport :
\lit{*} indique un Object vivant. \lit{\_} est le caractère de « kill »
d'Object.

La position peut être au format lat/long ou lat/long compressé, et le
rapport peut aussi contenir des données étendues. Les Objects compressés
ne sont pas recommandés sur RF en raison d'incompatibilités. Pour une
précision de 1 pied, utiliser le format \lit{!DAO!}. Un Object a toujours
un timestamp. Le champ Commentaire peut contenir n'importe quelle donnée
APRS appropriée (voir la section Champ Commentaire du chapitre 5 : Données
APRS dans le champ Information AX.25).

(Pourquoi le champ Commentaire est-il limité à 36 ou 43 caractères ? Le
chapitre 3 indique que la partie Information peut comporter jusqu'à 256
caractères.)

APRS 1.1 : la version compressée n'est pas recommandée.
\end{aprssideblock}

\begingroup
\setlength{\arrayrulewidth}{0.6pt}
\setlength{\tabcolsep}{2pt}
\renewcommand{\arraystretch}{1.5}
\sffamily\fontsize{7.5}{8.5}\selectfont
\begin{center}
\begin{tabular}{|>{\centering\arraybackslash}m{0.05\textwidth}|>{\centering\arraybackslash}m{0.11\textwidth}|>{\centering\arraybackslash}m{0.06\textwidth}|>{\centering\arraybackslash}m{0.10\textwidth}|>{\centering\arraybackslash}m{0.08\textwidth}|>{\centering\arraybackslash}m{0.08\textwidth}|>{\centering\arraybackslash}m{0.08\textwidth}|>{\centering\arraybackslash}m{0.07\textwidth}|>{\centering\arraybackslash}m{0.14\textwidth}|>{\centering\arraybackslash}m{0.13\textwidth}|}
\hline
\multicolumn{10}{|>{\columncolor{gray!20}}l|}{\hspace{0.6em}\bfseries\itshape Format Object Report --- avec position Lat/Long}\\
\hline
\chElevenHead{\lit{;}} & \chElevenHead{Object\\Name} & \chElevenHead{\lit{*} ou\\\lit{\_}} & \chElevenHead{Time\\DHM /\\HMS} & \chElevenHead{Lat} & \chElevenHead{Sym\\Table\\ID} & \chElevenHead{Long} & \chElevenHead{Symbol\\Code} & \chElevenDataExtHead & \cellcolor{gray!20}\chElevenHead{Comment\\(max 36 avec ext,\\43 sans)} \\
\hline
\chElevenBytes{1} & 9 & 1 & 7 & 8 & 1 & 9 & 1 & 7 & 0--36/43 \\
\hline
\multicolumn{10}{|p{0.96\linewidth}|}{%
  \vspace{0.25em}
  {\rmfamily\underline{Exemples}}

  \vspace{0.25em}
  \begin{tabular}{@{}>{\raggedright\arraybackslash}p{0.62\linewidth}>{\raggedright\arraybackslash}p{0.34\linewidth}@{}}
  \aprsorigtexttt{;LEADER   *092345z4903.50N/07201.75W>088/036} &
  \chElevenExampleText{un Object vivant. À 2345 heures zulu le 9 du mois, le
  « Leader » était dans la voiture à 49°3'30''N/72°1'45''W, cap 88 degrés à
  36 knots.} \\
  \aprsorigtexttt{;LEADER   \_092345z4903.50N/07201.75W>088/036} &
  \chElevenExampleText{le même Object, maintenant tué.} \\
  \end{tabular}
  \vspace{0.25em}
} \\
\hline
\end{tabular}
\end{center}
\endgroup

\begingroup
\setlength{\arrayrulewidth}{0.6pt}
\setlength{\tabcolsep}{2pt}
\renewcommand{\arraystretch}{1.5}
\sffamily\fontsize{7.5}{8.5}\selectfont
\begin{center}
\begin{tabular}{|>{\centering\arraybackslash}m{0.05\textwidth}|>{\centering\arraybackslash}m{0.13\textwidth}|>{\centering\arraybackslash}m{0.07\textwidth}|>{\centering\arraybackslash}m{0.10\textwidth}|>{\centering\arraybackslash}m{0.20\textwidth}|>{\centering\arraybackslash}m{0.30\textwidth}|}
\hline
\multicolumn{6}{|>{\columncolor{gray!20}}l|}{\hspace{0.6em}\bfseries\itshape Format Object Report --- avec position Lat/Long compressée --- non recommandé}\\
\hline
\chElevenHead{\lit{;}} & \chElevenHead{Object\\Name} & \chElevenHead{\lit{*} ou\\\lit{\_}} & \chElevenHead{Time\\DHM /\\HMS} & \chElevenHead{Compressed Position\\Data} & \cellcolor{gray!20}\chElevenHead{Comment} \\
\hline
\chElevenBytes{1} & 9 & 1 & 7 & 13 & 0--36/43 \\
\hline
\multicolumn{6}{|p{0.96\linewidth}|}{%
  \vspace{0.25em}
  {\rmfamily\underline{Exemple}}

  \vspace{0.25em}
  \begin{tabular}{@{}>{\raggedright\arraybackslash}p{0.62\linewidth}>{\raggedright\arraybackslash}p{0.34\linewidth}@{}}
  \aprsorigtexttt{;LEADER   *092345z/5L!!<*e7>7P[} &
  \chElevenExampleText{le « Leader » était dans la voiture à 49°30'00''N/72°45'00''W, cap 88 degrés à 36,2 knots.} \\
  \end{tabular}
  \vspace{0.25em}
} \\
\hline
\end{tabular}
\end{center}
\endgroup

\begin{aprssideblock}{Format Item Report\\(non recommandé)}
Un Item Report possède un nom d'Item de longueur variable, de 3 à
9 caractères. Le nom peut être constitué de n'importe quels caractères
ASCII imprimables, y compris des espaces incorporés, sauf \lit{!} ou
\lit{\_} car ils sont les terminateurs de champ. Les noms d'Item sont
sensibles à la casse. Noter que le champ de nom est de longueur variable,
contrairement aux objects, où il est fixe.

Le \lit{)} est l'APRS Data Type Identifier d'un Item Report, et un
\lit{!} ou un \lit{\_} sépare le nom d'Item du reste du rapport :
\lit{!} indique un Item vivant. \lit{\_} est le caractère de « kill »
d'Item.

La position peut être au format lat/long ou lat/long compressé. Il n'y a
pas de prévision de timestamp. Le rapport peut aussi contenir des données
étendues. Le champ Commentaire peut contenir n'importe quelle donnée APRS
appropriée (voir la section Champ Commentaire du chapitre 5 : Données APRS
dans le champ Information AX.25).

APRS 1.1 : le format Item n'est pas recommandé sur RF en raison
d'incompatibilités.
\end{aprssideblock}

\begingroup
\setlength{\arrayrulewidth}{0.6pt}
\setlength{\tabcolsep}{2pt}
\renewcommand{\arraystretch}{1.5}
\sffamily\fontsize{7.5}{8.5}\selectfont
\begin{center}
\begin{tabular}{|>{\centering\arraybackslash}m{0.05\textwidth}|>{\centering\arraybackslash}m{0.11\textwidth}|>{\centering\arraybackslash}m{0.06\textwidth}|>{\centering\arraybackslash}m{0.08\textwidth}|>{\centering\arraybackslash}m{0.08\textwidth}|>{\centering\arraybackslash}m{0.08\textwidth}|>{\centering\arraybackslash}m{0.07\textwidth}|>{\centering\arraybackslash}m{0.14\textwidth}|>{\centering\arraybackslash}m{0.13\textwidth}|}
\hline
\multicolumn{9}{|>{\columncolor{gray!20}}l|}{\hspace{0.6em}\bfseries\itshape Format Item Report --- avec position Lat/Long --- non recommandé}\\
\hline
\chElevenHead{\lit{)}} & \chElevenHead{Item\\Name} & \chElevenHead{\lit{!} ou\\\lit{\_}} & \chElevenHead{Lat} & \chElevenHead{Sym\\Table\\ID} & \chElevenHead{Long} & \chElevenHead{Symbol\\Code} & \chElevenDataExtHead & \cellcolor{gray!20}\chElevenHead{Comment\\(max 36 avec ext,\\43 sans)} \\
\hline
\chElevenBytes{1} & 3--9 & 1 & 8 & 1 & 9 & 1 & 7 & 0--36/43 \\
\hline
\multicolumn{9}{|p{0.96\linewidth}|}{%
  \vspace{0.25em}
  {\rmfamily\underline{Exemples}}

  \vspace{0.25em}
  \begin{tabular}{@{}>{\raggedright\arraybackslash}p{0.55\linewidth}>{\raggedright\arraybackslash}p{0.41\linewidth}@{}}
  \aprsorigtexttt{)AID\#2!4903.50N/07201.75WA} &
  \chElevenExampleText{la First Aid Station \#2 est à 49°3'30''N/72°1'45''W.
  (\lit{/A} est le symbole de Aid Station).} \\
  \aprsorigtexttt{)G/WB4APR!53VV.VVN\textbackslash002VV.VVWd} &
  \chElevenExampleText{une station DX rare « quelque part en Angleterre ».
  (\lit{\textbackslash d} est le symbole de DX Spot).} \\
  \aprsorigtexttt{)AID\#2\_4903.50N/07201.75WA} &
  \chElevenExampleText{la First Aid Station a fermé.} \\
  \end{tabular}
  \vspace{0.25em}
} \\
\hline
\end{tabular}
\end{center}
\endgroup

\begingroup
\setlength{\arrayrulewidth}{0.6pt}
\setlength{\tabcolsep}{2pt}
\renewcommand{\arraystretch}{1.5}
\sffamily\fontsize{7.5}{8.5}\selectfont
\begin{center}
\begin{tabular}{|>{\centering\arraybackslash}m{0.05\textwidth}|>{\centering\arraybackslash}m{0.13\textwidth}|>{\centering\arraybackslash}m{0.07\textwidth}|>{\centering\arraybackslash}m{0.20\textwidth}|>{\centering\arraybackslash}m{0.30\textwidth}|}
\hline
\multicolumn{5}{|>{\columncolor{gray!20}}l|}{\hspace{0.6em}\bfseries\itshape Format Item Report --- avec position Lat/Long compressée --- non recommandé}\\
\hline
\chElevenHead{\lit{)}} & \chElevenHead{Item\\Name} & \chElevenHead{\lit{!} ou\\\lit{\_}} & \chElevenHead{Compressed Position\\Data} & \cellcolor{gray!20}\chElevenHead{Comment} \\
\hline
\chElevenBytes{1} & 3--9 & 1 & 13 & 0--36/43 \\
\hline
\multicolumn{5}{|p{0.96\linewidth}|}{%
  \vspace{0.25em}
  {\rmfamily\underline{Exemple}}

  \vspace{0.25em}
  \begin{tabular}{@{}>{\raggedright\arraybackslash}p{0.55\linewidth}>{\raggedright\arraybackslash}p{0.41\linewidth}@{}}
  \aprsorigtexttt{)MOBIL!\textbackslash5L!!<*e79 sT} &
  \chElevenExampleText{la Mobil Gas Station est à 49°30'00''N/72°45'00''W.
  (\lit{\textbackslash 9} est le symbole de Gas Station).} \\
  \end{tabular}
  \vspace{0.25em}
} \\
\hline
\end{tabular}
\end{center}
\endgroup

\begin{aprssideblock}{Area Objects}
En utilisant le symbole \lit{\textbackslash l} (c'est-à-dire le symbole de la
lettre « L » minuscule de l'Alternate Symbol Table), il est possible de
définir des objects cercle, ligne, ellipse, triangle et boîte dans toutes
les couleurs, ouverts ou remplis, de toute taille allant de 60 pieds à
100 miles. Ces Objects sont utiles pour les événements en temps réel, comme
une recherche et sauvetage, ou l'ajout d'une route ou d'un itinéraire
spécial pour un événement spécial.

Le format d'Object est spécifié comme une APRS Data Extension de
7 caractères \texttt{Tyy/Cxx} placée immédiatement après le Symbol Code
\lit{l}. Par exemple :
\aprsorigtexttt{;OBJECT   *ddmm.hhN\textbackslash dddmm.hhW l Tyy/Cxx}

où :
\texttt{T} est le type de forme de l'object. \texttt{/C} est la couleur de
l'object. \texttt{yy} est la racine carrée de (l'offset de latitude en
degrés / 1500). \texttt{xx} est la racine carrée de (l'offset de longitude
en degrés / 1500).

À la réception :
offset\_en\_degrés = (offset\_dans\_le\_paquet\textsuperscript{2}) / 1500

\textbf{NOTE :} cette spécification avait à l'origine un facteur d'échelle
de 100. \url{http://www.aprs.org/aprs11/areaobjects.txt} contient la
correction pour le porter à 1500. Toutes les applications modernes
vérifiées utilisent 1500.
\end{aprssideblock}

\begingroup
\setlength{\arrayrulewidth}{0.6pt}
\setlength{\tabcolsep}{4pt}
\renewcommand{\arraystretch}{1.4}
\sffamily\fontsize{8.5}{9.5}\selectfont
\begin{center}
\begin{tabular}{|>{\centering\arraybackslash}m{0.12\textwidth}|>{\raggedright\arraybackslash}p{0.32\textwidth}|>{\centering\arraybackslash}m{0.12\textwidth}|>{\raggedright\arraybackslash}p{0.32\textwidth}|}
\hline
\rowcolor{gray!20}\bfseries T & \bfseries Type d'object & \bfseries T & \bfseries Type d'object \\
\hline
0 & Open circle & 5 & Color-filled circle \\
\hline
1 & Line (offset : down/right) & 6 & Line (offset : down/left) \\
\hline
2 & Open ellipse & 7 & Color-filled ellipse \\
\hline
3 & Open triangle & 8 & Color-filled triangle \\
\hline
4 & Open box & 9 & Color-filled box \\
\hline
\end{tabular}

\medskip

\begin{tabular}{|>{\centering\arraybackslash}m{0.18\textwidth}|>{\raggedright\arraybackslash}p{0.30\textwidth}|>{\centering\arraybackslash}m{0.18\textwidth}|}
\hline
\rowcolor{gray!20}\bfseries /C & \bfseries Couleur & \bfseries Intensité \\
\hline
/0 & Black & High \\
\hline
/1 & Blue & High \\
\hline
/2 & Green & High \\
\hline
/3 & Cyan & High \\
\hline
/4 & Red & High \\
\hline
/5 & Violet & High \\
\hline
/6 & Yellow & High \\
\hline
/7 & Gray & High \\
\hline
/8 & Black & Low \\
\hline
/9 & Blue & Low \\
\hline
/10 & Green & Low \\
\hline
/11 & Cyan & Low \\
\hline
/12 & Red & Low \\
\hline
/13 & Violet & Low \\
\hline
/14 & Yellow & Low \\
\hline
/15 & Gray & Low \\
\hline
\end{tabular}
\end{center}
\endgroup

\begin{aprsbodyblock}
\begin{itemize}\tightlist
\item
  Les triangles sont toujours des triangles isocèles, orientés
  verticalement, et les points de référence sont le sommet supérieur et le
  coin inférieur droit.
\item
  Les cercles sont définis par une boîte carrée référencée par le coin
  supérieur gauche et le centre.
\item
  Les lignes sont référencées par leur extrémité supérieure.
\end{itemize}

La position de « référence d'offset » de l'object est le coin supérieur
gauche de l'object, et les offsets sont la distance depuis le coin
inférieur droit (ou le centre d'un cercle) jusqu'à cette position de
« référence d'offset ». (Une exception est le cas particulier de la ligne
de type 6, qui est tracée vers le bas et vers la gauche.)

Voici quelques exemples de rapports de position d'Object. Les offsets de
latitude et de longitude valent chacun 4 minutes = 4/60 degré = 100/1500 de
degré), donc yy = xx = √100 = 10.
\end{aprsbodyblock}

\begin{aprsexample}
\aprsorigtexttt{;SEARCH   *092345z4903.50N\textbackslash 07201.75W l 710/310}

\chElevenExampleText{une ellipse remplie cyan haute intensité, yy=10, xx=10}

\aprsorigtexttt{;SEARCH   *092345z4903.50N\textbackslash 07201.75W l 8101310}

\chElevenExampleText{un triangle rempli violet basse intensité, yy=10, xx=10}
\end{aprsexample}

\begin{aprssideblock}{Area Objects\\(suite)}
En outre, avec l'option ligne (Type 1 et Type 6), il est possible de
spécifier un « corridor » de chaque côté de la ligne centrale. La largeur
du corridor (en miles) de chaque côté de la ligne est spécifiée dans le
texte du commentaire, entre \texttt{\{\}}. Par exemple :
\aprsorigtexttt{;FLIGHTPTH*4903.50N\textbackslash 07201.75W l 610/310\{100\}}

une ligne cyan haute intensité, avec un corridor de 100 miles de chaque
côté.

\textbf{Note :} l'option de remplissage de couleur doit être utilisée avec
précaution, car un object rempli de couleur masquera les informations
affichées en dessous.
\end{aprssideblock}

\begin{aprssideblock}{Signpost Objects / Items}
Les Signpost Objects/Items (avec le symbole \lit{\textbackslash m})
s'affichent comme une boîte jaune portant une superposition de 1 à
3 caractères. La superposition est spécifiée en plaçant les 1 à
3 caractères entre accolades dans le champ commentaire. Ainsi, un panneau
avec \{55\} apparaîtra comme un panneau portant 55. Par exemple :
\aprsorigtexttt{)I91 3N!4903.50N\textbackslash 07201.75Wm\{55\}}

Cela a été conçu à l'origine pour afficher la vitesse du trafic devant des
appareils de mesure de vitesse, mais peut servir à toute fin.

\textbf{Recommandation d'implémentation :} les panneaux ne doivent afficher
aucun indicatif ni nom, et pour éviter l'encombrement, ne doivent être
affichés qu'à courte portée.
\end{aprssideblock}

\begin{aprssideblock}{Format Object obsolète}
Certaines stations transmettent des rapports d'Object sans l'APRS Data
Type Identifier \lit{;}. Ce format est obsolète. Certains logiciels
peuvent encore décoder de telles données comme un Object, mais elles
doivent désormais être interprétées comme un Status Report.
\end{aprssideblock}
